\documentclass[11pt]{article}
\usepackage[margin=1in]{geometry}
\usepackage[T1]{fontenc}
\usepackage{lmodern,amsmath,amssymb,amsthm,mathtools,hyperref}
\setlength{\emergencystretch}{3em}
\newtheorem{theorem}{Theorem}[section]
\newtheorem{lemma}[theorem]{Lemma}
\newtheorem{corollary}[theorem]{Corollary}
\newtheorem{proposition}[theorem]{Proposition}
\newcommand{\E}{\mathbb E}
\newcommand{\tr}{\operatorname{tr}}
\title{Constructive code escape, exact moment repair,\\and arithmetic obstructions}
\author{Joshua Mark\thanks{ORCID: \href{https://orcid.org/0009-0003-5905-2382}{0009-0003-5905-2382}. AI tools assisted with derivation, independent review, computation and exposition.}}
\date{20 November 2026 --- research draft; priority not established}
\begin{document}\maketitle

\begin{abstract}
We investigate construction of the near-extremal sign distributions
needed for Hadamard extraction. We give exact Fourier moment repair
with quantified concentration cost, and prove that even its minimal
affine buffer fails on nonisotropic uniform nearest-order candidates
at sufficiently large orders. An additive-cover theorem rules out
arbitrary reweighting of the smallest binary linear seed as a way
around its power-of-two support barrier. Conversely, explicit
character rewiring escapes all traps in a special infinite family,
with an optimal-overlap straight bridge and constructive completion.
Odd-integer lattice witnesses provide a separate repair diagnostic
whose parity and Euclidean-shortness limitations are made explicit.
No general construction of the required distribution at unknown
orders is obtained, and the full Hadamard conjecture remains open.
Novelty is not asserted pending further prior-art review.
\end{abstract}
The standing extraction target is the following. For $N\ge8$ divisible
by four, an unbiased pairwise-independent law on $N-4$ sign bits
with collision at least $1/N-1/(1024N^4)$ constructs a Hadamard
matrix of order $N$. The proof is preserved in
\path{../round6/collision_stability.tex}. This pass tests actual
construction and correction mechanisms against that target, rather
than assuming its input exists.

\section{Exact correction of low-order moments}

Let $P$ be any probability law on the $d$-dimensional sign cube,
$M=2^d$, and $\chi_S(x)=\prod_{i\in S}x_i$. Write
\[
 m_S=\sum_xP(x)\chi_S(x),\qquad
 f(x)=\sum_{1\le|S|\le2}m_S\chi_S(x),\quad
 B=\max_x f(x),\quad L=\sum_{1\le|S|\le2}|m_S|.
\]
All quantities are computable from a finite law. In particular
$0\le B\le L$. Let $C(P)=\sum_xP(x)^2$.

\begin{theorem}[Nonnegative exact correction]
The law
\[
 Q(x)=\frac{P(x)+(B-f(x))/M}{1+B}
\]
is nonnegative, normalized, unbiased and pairwise independent. Its
total variation distance from $P$ is at most $B/(1+B)$, and
\begin{equation}\label{correction:collision}
 C(Q)=\frac1M+
 \frac{C(P)-(1+\sum_{1\le|S|\le2}m_S^2)/M}{(1+B)^2}
 \ \ge\ \frac{C(P)}{(1+B)^2}.
\end{equation}
\end{theorem}
\begin{proof}
Nonnegativity is immediate from $B\ge f$. Uniform cube orthogonality
gives $\sum f=0$ and $M^{-1}\sum f\chi_S=m_S$ for the specified
sets $S$, so the low-order moments cancel exactly. If $B>0$, define
$R(x)=(B-f(x))/(MB)$, a probability law. Then
$Q=(P+BR)/(1+B)$, proving the variation bound and pointwise
$Q\ge P/(1+B)$, hence the collision lower bound. If $B=0$, then
$f=0$ and $Q=P$. For the exact identity expand the square, using
$\sum P f=\sum m_S^2$ and $M^{-1}\sum f^2=\sum m_S^2$.
\end{proof}

The smallest nonnegative buffer in this same affine ansatz is exactly
\[
 \beta_* =\max\{0,\max_x(f(x)-MP(x))\}.
\]
Replacing $B$ by any $\beta\ge\beta_*$ preserves normalization,
exact low-order moment cancellation, and the equality in
\eqref{correction:collision}. The mixture interpretation and its simple
variation/collision lower bounds require $\beta\ge B$; they must not
be asserted for the potentially smaller $\beta_*$. Thus the explicit
universal construction is conservative, not an optimal repair theorem.

\begin{theorem}[A sufficient input window for exact extraction]
Let $d=n-1$, $n\equiv1\pmod4$, $n\ge5$, and $N=n+3$.
If $C(P)\ge1/N$ and $B\le1/(2048N^3)$, the corrected law satisfies
\[
 0\le1/N-C(Q)\le1/(1024N^4).
\]
Consequently the quantitative residue-one extraction theorem of round6
constructs a Hadamard matrix of order $N$ from $Q$.
It suffices to replace the assumption on $B$ by the stronger,
easier-to-check condition $L\le1/(2048N^3)$.
\end{theorem}
\begin{proof}
Exact isotropy of the augmented rows $(1,x)$ follows from the
corrected moments, and the residue collision bound gives $C(Q)\le1/N$.
The first theorem gives
\[
 1/N-C(Q)\le\frac1N\left(1-\frac1{(1+B)^2}\right)
 \le 2B/N\le1/(1024N^4).
\]
Apply the round6 stability theorem. This implication assumes the two
stated input properties and does not produce such an input in unknown
Hadamard orders.
\end{proof}

\begin{theorem}[Quantization obstruction for a uniform row candidate]
In the same parameters, suppose $P$ is uniform on $N$ distinct sign
rows, and at least one first or second moment is nonzero. The universal
$B$-buffer correction above satisfies
\[
 1/N-C(Q)\ge1/(2N^2),
\]
so it cannot enter the round6 extraction window.
\end{theorem}
\begin{proof}
Since $N$ is even, each moment is an even integer divided by $N$;
a nonzero one has magnitude at least $2/N$. Uniform expectation
gives $|m_S|=|\mathbb E_{\rm unif} f\chi_S|
\le\mathbb E_{\rm unif}|f|\le2B$, since $f$ has mean zero
and positive part at most $B$. Thus $B\ge1/N$.
Here $C(P)=1/N$ and $M=2^{N-4}\ge2N$ for $N\ge8$.
The exact collision identity yields
\[
 1/N-C(Q)\ge(1/N-1/M)\left(1-(1+B)^{-2}\right)
 \ge\frac1{2N}\left(1-(1+1/N)^{-2}\right)
 \ge\frac1{2N^2}.
\]
The last inequality is equivalent to $N^2-N-1\ge0$.
\end{proof}
This obstruction concerns the specified guaranteed buffer, not every
correction scheme or the smaller $\beta_*$. It explains why repairing
an arbitrary nearly orthogonal nearest-order sign matrix by adding a
uniform Fourier buffer is inadequate. For uniform $N$-row candidates,
the sufficient error condition of the preceding theorem forces all
low-order moments to vanish already by integer quantization.
\section{Why even the minimal affine Fourier buffer is insufficient}
\label{sec:minimal-buffer-obstruction}
The preceding Fourier correction is useful for appropriately weighted,
very accurate inputs. Here we show a limitation that also applies to
the \emph{smallest} nonnegative buffer, rather than just its universal
upper bound. The argument uses elementary low-degree concentration;
it does not exclude other nonlinear correction procedures.

Let $N\ge20$ be a multiple of four, $d=N-4$, $M=2^d$ and
$K=d(d+1)/2$. Let $P$ be uniform on $N$ distinct points of the
$d$-bit cube. For $1\le|S|\le2$, let
$m_S=\mathbb E_P\chi_S$, and put
\[
 f(x)=\sum_{1\le|S|\le2}m_S\chi_S(x),\qquad
 \sigma^2=\sum_Sm_S^2.
\]
For any feasible buffer $\beta\ge0$, consider
\[
 Q_\beta(x)=\frac{P(x)+(\beta-f(x))/M}{1+\beta}.
\]
Its least nonnegative choice is
$\beta_* =\max(0,\max_x(f(x)-MP(x)))$.

\begin{theorem}[Sparse uniform candidates cannot enter the extraction window]
\label{thm:minimal-buffer-obstruction}
If $P$ is not already unbiased and pairwise independent, then every
nonnegative correction $Q_\beta$ in this family satisfies
\[
 \frac1N-C(Q_\beta)\ge\frac1{8N^2\sqrt K}>
 \frac1{8N^3}>\frac1{1024N^4}.
\]
In particular none enters the sufficient window proved in round six.
If $P$ is already exactly isotropic, its uniform $N$ atoms already
give a Hadamard matrix by the residue equality theorem.
\end{theorem}
\begin{proof}
All expectations in the next paragraph are under the uniform cube.
Fourier orthogonality gives $\mathbb E f=0$ and
$\mathbb E f^2=\sigma^2$. Cauchy--Schwarz gives
$\|f\|_\infty\le\sqrt K\sigma$, so
\[
 \mathbb E|f|\ge\frac{\sigma}{\sqrt K},\qquad
 \mathbb E f_+=\tfrac12\mathbb E|f|
       \ge\frac{\sigma}{2\sqrt K}.
\]
Let $T$ be the $N$-point support of $P$. Its contribution to
$\mathbb E f_+$ is at most $(N/M)\sqrt K\sigma$.
Since $M\ge4NK$, a positive contribution of at least
$\sigma/(4\sqrt K)$ remains outside $T$. Hence
\[
 \beta_*\ge\max_{x\notin T}f(x)
          \ge\frac{\sigma}{4\sqrt K}.
\]
A nonzero $m_S$ is an even integer divided by $N$; therefore
$\sigma\ge2/N$, and
$\beta\ge\beta_0:=1/(2N\sqrt K)$.

The exact collision identity for the correction, valid for every
feasible $\beta$, is
\[
 C(Q_\beta)=\frac1M+
 \frac{1/N-(1+\sigma^2)/M}{(1+\beta)^2}
 \le\frac1M+\frac{1/N-1/M}{(1+\beta_0)^2}.
\]
Here the negative $\sigma^2$ term was discarded \emph{before}
comparing denominators, so no unproved sign assumption is needed.
Since $M\ge2N$, $\beta_0\le1$, and
$1-(1+\beta_0)^{-2}\ge\beta_0/2$, subtraction from $1/N$ gives
the claimed lower bound. Also $\sqrt K<N$.

For completeness, $2^{16}=65536\ge4\cdot20\cdot136$ verifies
$M\ge4NK$ at $N=20$. Increasing $N$ by four multiplies $M$ by
$16$, while $NK$ grows by a factor at most
$(24/20)(20\cdot21)/(16\cdot17)=63/34<2$.
Thus the inequality holds for all stated $N$.
\end{proof}

The conclusion concerns sparse \emph{uniform} starting laws. It does
not negate the correction theorem for nonuniform weighted candidates
whose initial concentration and moment accuracy are both sufficient.
Finding such candidates is still an existence problem, not a
consequence of the correction identity.
\section{An additive-cover obstruction to small binary-seed constructions}
\label{sec:character-rigidity}
This section concerns a specific construction class, not arbitrary
pairwise-independent laws. The observed bits are fixed linear
characters of a binary seed; the seed distribution itself is allowed
to be completely nonuniform. This distinction matters because
rewiring the observed characters can escape a code trap, whereas
changing only the seed weights may be impossible.

Let $G=\mathbb F_2^r$, $L=2^r$, and let
$B\subset G\setminus\{0\}$ be $d$ distinct masks spanning $G$.
For a seed law $p$ write
\[
 \chi_b(z)=(-1)^{b\cdot z},\qquad
 \widehat p(c)=\sum_{z\in G}p_z\chi_c(z),\qquad
 A=\{0\}\cup B.
\]
The observed sign vector is $X(z)=(\chi_b(z))_{b\in B}$.
Spanning makes this map injective, so its collision probability equals
$\sum_zp_z^2$.

\begin{theorem}[Complete rigidity from additive coverage]
\label{thm:character-cover}
The observed bits are unbiased and pairwise independent if and only if
\[
 \widehat p(c)=0\quad\text{for every }c\in(A+A)\setminus\{0\}.
\]
Consequently, if $A+A=G$, the only feasible seed law is uniform.
In particular this holds whenever $|A|>L/2$.
If $c\notin A+A$, the family
\[
 p_t(z)=\frac{1+t\chi_c(z)}L,\qquad -1\le t\le1,
\]
is feasible and has collision $(1+t^2)/L$. At $t=\pm1$ it is uniform
on an affine hyperplane of size $L/2$ and doubles the collision.
Thus missing additive directions give explicit compression steps,
whereas full coverage forbids every seed-reweighting step.
\end{theorem}
\begin{proof}
The coordinate means are $\widehat p(b)$, and distinct pair
correlations are $\widehat p(b+b')$. These masks, together with zero,
are exactly $A+A$. If they cover $G$, Fourier inversion forces
$p_z=1/L$. If $|A|>L/2$, then for each $c\in G$ the sets
$A$ and $A+c$ intersect, implying $c\in A+A$.
For a missing $c$, character orthogonality gives the displayed law
normalization, nonnegativity, vanishing required Fourier coefficients,
and collision. Moreover $c\ne0$, so its two level sets are affine
hyperplanes of size $L/2$. No selected mask is zero or identified
with another after restriction: either failure would put $c$ in
$A+A$. Thus the compression retains all observed moments.
\end{proof}

The character construction and finite-group Fourier identities are
classical; see Luby--Wigderson, \emph{Pairwise Independence and
Derandomization}, Sections 1.1--1.2,
\url{https://www.math.ias.edu/~avi/PUBLICATIONS/LubyWig2005.pdf}.
The proof above is included to make the precise construction-class
restriction explicit; no priority claim is made.

\begin{corollary}[The smallest binary seed forces a power-of-two law]
\label{cor:minimum-seed}
If $r=\lceil\log_2(d+1)\rceil$, any feasible law with these linear
observables is uniform on $L=2^r$ distinct atoms, even when arbitrary
seed weights are allowed.
For a non-power-of-two multiple of four $N\ge12$, take $d=N-4$.
Then this construction class has collision deficit
\[
 \frac1N-C=\frac1N-\frac1L>\frac2{N^2},
\]
far larger than the sufficient extraction window $1/(1024N^4)$.
\end{corollary}
\begin{proof}
Minimality of $r$ implies $d+1>L/2$, so apply the theorem.
For the stated $N$, $L\ge N+4$ and $L<2N$, since both $L,N$ are
multiples of four and $N-3$ lies strictly between the preceding and
current powers of two. Hence
$(L-N)/(NL)\ge4/(NL)>2/N^2$.
\end{proof}
For example $N=12,d=8$ forces sixteen equally weighted atoms in the
smallest binary seed, with collision $1/16$ instead of $1/12$.
This happens even though an order-twelve Hadamard matrix exists.
The obstruction belongs to the linear-observable model; it is not
an obstruction to Hadamard existence.

\begin{proposition}[Approximate moment version]\label{prop:approx-character}
Suppose $A+A=G$ and every observed mean and pair correlation has
absolute value at most $\varepsilon$. Then
\[
 C(p)=\frac1L\left(1+\sum_{c\ne0}\widehat p(c)^2\right)
 \le\frac1L+\frac{L-1}{L}\varepsilon^2.
\]
In particular, if $C(p)\ge1/N-\eta$ with $N<L$, necessarily
\[
 \varepsilon^2\ge
 \frac{L/N-1-L\eta}{L-1}
\]
whenever the right side is positive.
\end{proposition}
\begin{proof}
Coverage makes each nonzero Fourier coefficient an observed mean or
pair correlation. Parseval identity gives the formula and the bound.
Rearranging gives the last assertion.
\end{proof}
Thus allowing very small moment errors does not silently bypass the
power-of-two barrier in this model. A broader distribution strategy
must change the observables, use a larger seed with nontrivial Fourier
freedom, or employ genuinely nonlinear constructions. The theorem
does not rule out any of those alternatives.
\section{Optimal-overlap escape from the character traps}
\label{sec:compression-escape}
This section gives an explicit nonlocal escape for the special character
traps of the preceding pass. It does not produce the missing distributions
at arbitrary admissible orders. The only background is orthogonality of
finite-group characters; a primary account of pairwise-independent finite
field constructions is Luby--Wigderson, Sections 1.1--1.2,
\url{https://www.math.ias.edu/~avi/PUBLICATIONS/LubyWig2005.pdf}.
The Walsh and block constructions here are classical ingredients; priority
for these landscape statements has not been established.

Put $h=2^s$, $s\ge2$, and identify the two seed spaces with
$V=\mathbb F_2^s$. Fix a nonzero linear functional $e$ on $V$.
For independent uniform $a,b\in V$, the trap law $\mu$ has observed bits
\[
 (\chi_\alpha(a))_{\alpha\ne0},\qquad
 (\chi_\beta(b))_{\beta\ne0},\qquad
 \chi_e(a+b),
 \quad\text{where }\chi_\alpha(a)=(-1)^{\alpha(a)}.
\]
There are $2h-1$ bits and $h^2$ distinct atoms. Append the constant
coordinate and write $r(a,b)$ for the augmented sign row of length $2h$.
The endpoint law $\nu$ is instead generated by uniform $a\in V$ and
$t\in\mathbb F_2$:
\begin{equation}\label{eq:compression-endpoint}
 (\chi_\alpha(a))_{\alpha\ne0},\qquad
 ((-1)^t\chi_\beta(a))_{\beta\ne0},\qquad (-1)^t.
\end{equation}
Its masks are exactly the nonzero characters on $V\oplus\mathbb F_2$,
so it is balanced pairwise independent and uniform on $2h$ atoms.
This directly constructs the compressed endpoint.

\begin{theorem}[Maximum overlap and optimal straight bridge]
\label{thm:compression-optimal-overlap}
Every pairwise-orthogonal subset of augmented trap rows has size at most
$h$, and $h$ is attained. Consequently every global collision maximizer
on $2h-1$ unbiased pairwise-independent bits has at most $h$ atoms in
common with the trap. The endpoint \eqref{eq:compression-endpoint}
attains this maximum overlap. Along the feasible affine bridge
$\mu_\lambda=(1-\lambda)\mu+\lambda\nu$,
\[
 C(\mu_\lambda)=\frac{1-\lambda+(h/2)\lambda^2}{h^2},\qquad
 \min_{0\le\lambda\le1}C(\mu_\lambda)
 =\frac{1-1/(2h)}{h^2},
\]
with the minimum at $\lambda=1/h$.
Among all straight mixture bridges from this trap to a global maximizer,
this minimum is as large as possible. The optimal collision loss is
exactly $1/(2h^3)$.
\end{theorem}
\begin{proof}
Character orthogonality gives
\[
 r(a,b)\cdot r(a',b')=
 h\mathbf1_{a=a'}+h\mathbf1_{b=b'}-1+
 (-1)^{e(a+a'+b+b')}.
\]
Distinct orthogonal rows therefore have different $a$ coordinates,
different $b$ coordinates, and the same value of $e(a+b)$.
In particular there are at most $h$. All diagonal seeds $(a,a)$ have
mixed bit $+1$ and form an orthogonal $h$-set.

The endpoint contains precisely these diagonal trap atoms in its $t=0$
half. Its $t=1$ half has no trap atoms: that would require
$\chi_\beta(b)=-\chi_\beta(a)$ for every nonzero $\beta$, impossible
when $s\ge2$, since multiplying the conditions for two distinct nonzero
functionals contradicts the condition for their sum. Thus the overlap
has exactly $h$ atoms and
$\mu\cdot\nu=h/(h^2\cdot2h)=1/(2h^2)$.
Expanding the squared norm proves the bridge formula and minimum.

The global collision bound is $1/(2h)$, attained by $\nu$. Indeed the
standard atom bound is $w_x\le1/(2h)$; equality in
$\sum w_x^2\le1/(2h)$ forces uniform weights on $2h$ atoms, and their
augmented rows must be orthogonal. Every other global endpoint therefore
has collision $1/(2h)$ and cross inner product with $\mu$ at most
$1/(2h^2)$. Its mixture collision is pointwise no larger than the
formula above. Evaluating at $\lambda=1/h$ proves optimality among
straight bridges. No optimality over arbitrary curved paths is asserted.
\end{proof}

\begin{proposition}[Complete any maximum overlap subset]
\label{prop:compression-complete-clique}
Every orthogonal $h$-subset of the augmented trap support extends to an
order-$2h$ Hadamard matrix with exactly $h$ trap rows.
\end{proposition}
\begin{proof}
The preceding characterization says its seed pairs form the graph of a
bijection $b=\pi(a)$ with constant mixed bit $(-1)^{e(a+\pi(a))}=c$.
Arrange the columns as $(1,A\text{-bits})$ and
$(\text{mixed bit},B\text{-bits})$. The selected rows are
$[H_1\ H_2]$, where $H_1$ is the Walsh matrix of $V$ and $H_2$ is
another such matrix with permuted rows and its constant column multiplied
by $c$. Both are Hadamard matrices of order $h$. Append the $h$ rows
$[H_1\ -H_2]$. Direct block multiplication gives an order-$2h$ Hadamard
matrix. It contains the selected trap rows and cannot contain more than
$h$ by the theorem.
\end{proof}

\subsection{Moment-preserving jumps with nondecreasing endpoint collision}
The same escape can be organized into elementary character substitutions.
This is a discontinuous algorithm on feasible distributions, not a
continuous collision-ascent flow. Label a character mask by a pair of
functionals on $V\oplus V$. Initially the observed masks are
\[
 (\alpha,0)\ (\alpha\ne0),\quad (0,\beta)\ (\beta\ne0),\quad(e,e).
\]
First swap the columns labeled $(0,e)$ and $(e,e)$. Then for each
$\beta\ne0,e$, replace the mask in its $B$-column by $(\beta,e)$.
All masks remain distinct and nonzero throughout. The final mask list is
$(\alpha,0)$, $(\beta,e)$ for $\beta\ne0$, and $(0,e)$, which generates
the endpoint law above.

\begin{proposition}[Finite nonlocal rewiring]\label{prop:compression-jumps}
The $h-1$ jumps just described preserve unbiased pairwise independence
at every endpoint. Endpoint collision never decreases. Exactly $s-1$
jumps halve the support size and double collision; the other $h-s$
jumps preserve support size. If each jump is interpolated by its straight
feasible segment, the entire piecewise affine path has minimum collision
$3/(4h^2)$.
\end{proposition}
\begin{proof}
A uniform linear-character law with mask-span rank $r$ is uniform on
$2^r$ atoms. After the initial swap, every replacement mask lies in the
span of the unchanged first-block masks and the unchanged mask $(0,e)$.
Thus no replacement increases rank; each drops it by at most one.
The initial and final ranks are $2s$ and $s+1$, proving the counts.

For a one-column replacement let $r_0$ be the rank of the retained masks.
The new mask is in their span. If the old mask is outside that span, the
new support is a codimension-one subcode of the old support of size $m$.
The cross inner product of the uniform laws is $1/m$, so the segment
collision is $(1+\lambda^2)/m$. If the old mask is inside the retained
span, both supports have size $m$; their intersection has size $m/2$,
since two distinct linear functionals disagree on half the retained
code. The segment collision is
$(1-\lambda+\lambda^2)/m$, minimized at $3/(4m)$.
For the initial swap the supports also have size $h^2$ and intersect in
$h^2/2$ atoms: the unchanged masks determine the seed, and the swap
fixes its row precisely when $\chi_e(a)=1$. (The nonzero $B$ masks
other than $e$ still span the second seed space for $s\ge2$.) Hence
its segment has minimum $3/(4h^2)$. Subsequent $m\le h^2$ cannot
produce a smaller minimum.
\end{proof}

\paragraph{What is and is not compressed.}
A linear quotient of the original seed that retains all observed masks
cannot do this: the observed masks and their pairwise sums cover the
entire dual group, so every nonzero quotient kernel identifies two masks
or kills a mask. The algorithm changes the observed characters instead.
It exploits the special two-block structure; no corresponding rewiring
rule for arbitrary admissible dimensions is proved. The direct bridge
has a much smaller collision loss than the column-by-column path, while
the jump endpoints themselves never lose collision.
\section{Parity witnesses, weighted repair budgets, and their exact limit}
\label{sec:parity-repair}

The odd-block argument extends to arbitrary integer witnesses, including
overlapping or unequal odd blocks. The extension supplies explicit
constraints on \emph{which} rows a repair must remove. Saturated-lattice
duality also identifies exactly what these certificates detect: absence
of an all-odd integer kernel vector. A certified order-twenty-four
example shows why that arithmetic test still falls short of sign
extension. Arithmetic completion is a classical subject; a relevant
primary source is Verheiden's thesis,
\emph{Arithmetical Properties of Combinatorial Matrices} (1978),
\url{https://thesis.caltech.edu/2899/}, and his article
\emph{Integral and rational completions of Hadamard matrices},
\url{https://doi.org/10.1016/0097-3165(80)90028-X}.
The results below have self-contained proofs; no novelty claim is made.

\begin{theorem}[Weighted odd-witness budget]
Let $P$ have $r$ orthogonal sign rows $p_1,\ldots,p_r$ of length $N$,
and let $f\in\mathbb Z^N$ have odd coordinate sum. Set
\[
 R=NI_N-P^{\mathsf T}P,\qquad
 C_f=f^{\mathsf T}Rf
     =N\|f\|^2-\sum_i(p_i\cdot f)^2\ge0.
\]
If rows indexed by $D$ are deleted and $L$ mutually orthogonal sign
rows are inserted, orthogonal also to all retained rows, then
\[
 L\le C_f+\sum_{i\in D}(p_i\cdot f)^2.
\]
Consequently full completion requires
\[
 \sum_{i\in D}\bigl((p_i\cdot f)^2-1\bigr)
       \ge N-r-C_f.
\]
Every odd-sum integer $f$ gives such an inequality, so a collection
of witnesses gives simultaneous weighted constraints on row selection.
\end{theorem}
\begin{proof}
Bessel's inequality for orthogonal rows of squared norm $N$ bounds
the sum of squared inner products with $f$ by $N\|f\|^2$.
For every sign row $v$, $v\cdot f\equiv\sum_j f_j\equiv1\pmod2$,
so $(v\cdot f)^2\ge1$. Subtracting the contributions of the retained
rows proves the first inequality. A full completion inserts
$L=N-r+|D|$, giving the second.
\end{proof}

For a signed indicator $f$ of an odd block of size $s$, the constant
is $C_f=Ns-\sum_i(p_i\cdot f)^2$. The theorem applies independently
to every block, with no equality-of-sizes or disjointness assumption.
If $f$ lies in the rational row span of $P$, then $C_f=0$ and no
additional sign row exists. The budget then depends only on the
deleted rows' squared overlaps with the witness.

\begin{proposition}[The first possible parity deficit is eight]
Assume $4\mid N$ and write $d=N-r$. For every odd-sum integer $f$,
\[
 C_f\equiv d\pmod8.
\]
Thus for $0\le d\le7$ these witness inequalities automatically give
$C_f\ge d$; they cannot certify failure of full completion without
row deletion. At $d=8$, the value $C_f=0$ is possible.
\end{proposition}
\begin{proof}
Every $p_i\cdot f$ is odd, hence its square is $1$ modulo $8$.
Also $\|f\|^2$ is odd, so $N\|f\|^2\equiv N\pmod8$ because
$4\mid N$. Therefore $C_f\equiv N-r\pmod8$.
Nonnegativity gives the remaining assertion. The repeated $4\times12$
seed supplies the value zero at defect eight. This congruence alone
does not prove the classical seven-row completion theorem.
\end{proof}

\subsection{What the certificate detects exactly}
\begin{theorem}[Saturated-lattice alternative]
For any rational row space $V\subseteq\mathbb Q^N$, exactly one of
the following holds:
\begin{enumerate}
\item $V\cap\mathbb Z^N$ contains a vector with odd coordinate sum;
\item $V^\perp\cap\mathbb Z^N$ contains a vector all of whose
coordinates are odd.
\end{enumerate}
For a partial Hadamard row space, the first case supplies a zero-budget
witness. The second case supplies arithmetic parity feasibility,
but does not assert a sign vector exists.
\end{theorem}
\begin{proof}
Put $\Lambda=V\cap\mathbb Z^N$ and
$K=V^\perp\cap\mathbb Z^N$. Both are primitive: an integer multiple
lying in the relevant rational subspace implies the original vector
lies there. A primitive subgroup of $\mathbb Z^N$ has a basis
extendible to an integer unimodular basis, so reduction modulo $2$
preserves its rank. If $\dim V=r$, the images $\bar\Lambda,\bar K$
in $\mathbb F_2^N$ have dimensions $r,N-r$ and are orthogonal.
Thus $\bar K=\bar\Lambda^\perp$. The all-ones vector lies in
$\bar K$ exactly when every vector of $\bar\Lambda$ has even
coordinate sum. This is precisely the stated alternative.
\end{proof}

There is a simple exact algorithm for this alternative. Starting with
an independent integer row basis, find a dependence modulo $2$.
The corresponding sum of basis rows has even entries; divide it by
$2$ and replace any participating row. Rational span is preserved,
and the squared covolume decreases by a factor $4$. As it stays a
positive integer, the procedure terminates with independent reductions
modulo $2$. Its row lattice has odd index in the saturated lattice,
so it has the same reduction modulo $2$. An odd-sum basis row is
therefore a complete certificate for the first case. Otherwise select
a full-rank odd minor, set free kernel coordinates to $1$, and solve
for the pivot coordinates. Their denominators are odd, and clearing
those denominators produces the all-odd integer kernel vector.
The accompanying checker implements this procedure from scratch.

\subsection{Sparse certificates beyond uniform repeated seeds}
For the earlier terminal $16\times24$ frame, encoded by row masks
\[
\begin{gathered}
0,16773120,16519104,10359454,7684730,5820908,9389414,1547504,\\
6736278,11838780,3701846,13790378,10646512,2963882,12892774,5974836,
\end{gathered}
\]
with bit $j$ indicating entry $-1$ at zero-based coordinate $j$,
there are two particularly short witnesses:
\[
 f=e_2+e_8+e_{19},\qquad g=e_9-e_{10}+e_{22}.
\]
Direct integer calculation gives $C_f=C_g=0$ and
$\|f\|^2=\|g\|^2=3$. This certifies unextendibility without
enumerating candidate sign rows. The row indices having absolute
overlap $3$ are respectively
\[
 A=\{0,4,5,6,7,11,15\},\qquad
 B=\{4,6,7,9,10,14,15\};
\]
all other overlaps have magnitude $1$. Any augmenting exchange must
delete a row from each of $A$ and $B$. These necessary conditions
do not by themselves establish the previously certified three-row
exchange barrier, but they give concrete filters for a repair search.

There is also a positive case where the weighted budget selects an
optimal constructive route. In the repeated order-twelve tensor
family after $k$ fibers have been repaired, a triple-indicator
witness has zero budget. Original repeated rows have overlap
magnitude $3$, while installed affine-plane rows have overlap
magnitude $1$. Deleting $a$ original rows and $b$ installed rows
can therefore insert at most $9a+b$ rows. Net growth is at most
$8a$, so deleting only installed rows cannot augment the frame.
Full completion requires at least $t-k$ more original-row deletions,
regardless of how many installed rows are removed. The explicit fiber
repair attains this bound: sacrifice one original row and insert the
nine words $(x,y,x+y,x-y)$ over $\mathbb F_3$, interpreted as one
negative position per triple. This supplies a certified choice of
which class of rows to preserve during constructive repair.

\subsection{Odd integer feasibility still misses a Euclidean gap}
The lattice alternative has an explicit limitation. Consider the
$12\times24$ partial frame with row masks
\[
\begin{gathered}
0,16773120,16519104,9656012,13838682,14783914,\\
1944718,11409518,15230494,5921970,3487346,13018356.
\end{gathered}
\]
Its rows are orthogonal. Its integer kernel contains
\[
\begin{split}
v={}&(-1,-1,3,-1,1,1,-3,-1,1,1,-1,1,\\
&1,-1,-1,-1,1,-1,1,-1,1,-1,1,1),
\end{split}
\]
with $\|v\|^2=40$. Nevertheless its odd integer kernel has no vector
of squared norm $24$ or $32$, so its shortest odd vector has squared
norm exactly $40$.

Here this last assertion is a finite exact certificate. Rational
elimination selects twelve free coordinates. Enumerating their
$2^{11}=2048$ sign assignments up to global negation rules out norm
$24$. A norm-$32$ all-odd vector would have exactly one coordinate
of magnitude $3$ and every other coordinate of magnitude $1$.
If that exceptional coordinate is a pivot, it is covered by the same
2048 assignments. If it is free, there are another
$12\cdot2^{11}=24576$ assignments to check. None gives norm $32$.
Every all-odd vector has squared norm congruent to $24$ modulo $8$,
and the displayed vector attains the next possibility, $40$.
The verifier performs all these checks with exact integer and rational
arithmetic. Thus the frame has no sign extension despite passing the
all-odd kernel feasibility test. Real progress beyond parity requires
controlling the short-vector problem, not merely clearing denominator
or congruence obstructions.

\section{The remaining construction problem}
The special-family escape is a genuine construction and an optimal
straight mixture bridge, but its endpoint has power-of-two order.
Fixed linear observables in the smallest binary seed force uniformity,
and the minimal-buffer theorem prevents a uniform nonisotropic
nearest-order row candidate from reaching the extraction window by
that correction alone. The positive correction theorem instead
requires concentrated, accurately weighted inputs; their general
construction is still missing.

Arithmetic witnesses can diagnose and guide some row exchanges,
but absence of a parity obstruction need not give a sign row: a
shortest-vector issue remains in the odd kernel coset. These results
suggest testing larger or nonlinear seeds, adaptive corrections
beyond the one-buffer ansatz, and metric-aware lattice repairs.
They do not prove that any of these strategies works in every
admissible dimension. No Lean code or Git commit was made.
\end{document}
